\sscp{\tsc{Havu-Dhao} and Bima}{03-03-04}
Before we conclude this chapter we need to point
out a possible connection between \tsc{Havu-Dhao} and Bima.
There are two (possibly four) changes that are
shared between Bima and \tsc{Havu-Dhao},
but not shared with \tsc{Sumba} or other \tsc{East Sumbawa} languages.
How these shared changes came about remains to be investigated.
It could be due to parallel development
(the position taken by \citealp[93f]{Blust2008}),
or it could be due to shared development, strengthened by contact with Bima,
given that Bima historically had
the role of a dominant language.

The changes potentially shared between \tsc{Havu-Dhao} and Bima are
*t > **d, \it{t} (Dhao **d > \it{ɖ͡ʐ}); *r > \it{r,
l}; *p > **f, \it{p} (\tsc{Havu-Dhao} **f > {\0});
and *s > \it{ʧ}, \it{s}.
Each of these changes involve splits,
and for all changes there are words which do not have the same
reflexes in Bima vs. \tsc{Havu-Dhao}.
The different reflexes are summarised in \trf{03-15},
along with the number of reflexes in cognate sets for which a
reflex has been identified in both Bima and \tsc{Havu-Dhao}.

\begin{table}
	\centering\tcp{Bima and \tsc{Havu-Dhao} *l, *t, *p, and *s}{03-15}
	\st{5pt}\begin{tabular}{lll|lll|lll|lll}\lsptoprule
Bima&H-D&no.&Bima&H-D&no.&Bima&H-D&no.&Bima&H-D&no.\\
*l&*l&&*t&*t&&*p&*p&&*s&*s&\\	\midrule
{\hr}l&{\hr}l&11&{\hr}t&{\hr}t&3&{\hr}p&{\hr}p&10&{\hr}s&{\hr}h - s&4\\
{\hr}l&{\hr}r&1&{\hr}t&{\hr}d&3&{\hr}p&{\hr}{\0}&2&{\hr}s&{\hr}h - ʧ&1\\
{\hr}r&{\hr}r&12&{\hr}d&{\hr}d&18&{\hr}f&{\hr}{\0}&7&{\hr}ʧ&{\hr}h - ʧ&5\\
{\hr}r&{\hr}l&4&{\hr}d&{\hr}t&1&{\hr}f&{\hr}p&1&{\hr}ʧ&{\hr}h - s&3\\
		\lspbottomrule
	\end{tabular}
\end{table}

Positing shared development between Bima
and \tsc{Havu-Dhao} for *p requires also positing *p > **f > {\0} in \tsc{Havu-Dhao}.
While this is a common sequence of changes,
there is no independent evidence that it occurred in \tsc{Havu-Dhao}.
Similarly, the change of *s > \it{ʧ} is only shared between Bima
and Dhao, as Havu has undergone universal *s > \it{h}.
While some instances of *s > \it{h} may have been **ʧ at an earlier
stage, there is no independent evidence that this was the case.

Leaving aside the reflexes of *s,
while Bima and \tsc{Havu-Dhao} do not always have the same reflexes,
the disagreements between Bima
and \tsc{Havu-Dhao} are not random.\footnote{
		The discussion in \citet[93f]{Blust2008}
		is focussed on *p.
		%(Blust only mentions the changes affecting *l
		%and *t in passing).
		He concludes that, once conditioning environments are accounted
		for, the number of cases of agreement are insufficient to establish
		that the changes are not due to chance.
		However, if we leave aside shared retentions -- which
		would not be due to shared development -- the
		%it should be clear from \trf{03-15}
		number of cases of innovation in \trf{03-15}
		provide a basis for positing some kind of shared development.}
Whenever \tsc{Havu-Dhao} *l has changed,
it has also changed in Bima,
while there is only one case in which \tsc{Havu-Dhao} has *l
> \it{r} but Bima retains *l = \it{l}.
Similarly, whenever *t and *p have changed in Bima,
they have also changed in \tsc{Havu-Dhao},
while there is only one case each in which Bima is innovative
but \tsc{Havu-Dhao} is conservative.
This means that we can posit a shared period of history between
Bima and \tsc{Havu-Dhao} during which the changes *l > \it{r},
*t > \it{d}, and *p > **f occurred.
After this shared period of history,
Bima had additional cases of *l > \it{r},
while \tsc{Havu-Dhao} had additional cases of *t > \it{d}
and *p > **f > {\0}.

Given the proposed subgrouping in this chapter,
%links between Bima and \tsc{East Sumbawa}
%on one hand, and the links between \tsc{Havu-Dhao} and \tsc{Sumba}
%on the other hand,
the putative period of history shared
between \tsc{Havu-Dhao} and Bima was probably a period
of shared contact these languages had.
This contact would have taken place after
Bima had split from Proto-East Sumbawa
and \tsc{Havu-Dhao} had split from Proto-Sumba-Havu.\footnote{
	Alternately, we could propose
	Bima and \tsc{Havu-Dhao} subgroup together
	to the exclusion of \tsc{East Sumbawa} and \tsc{Sumba} respectively.
	This would then entail shared contact between \tsc{Havu-Dhao}
	and \tsc{Sumba} on the one hand, and Bima
	and \tsc{East Sumbawa} on the other hand.}

Examples of shared *t > **d (> \it{ɖ͡ʐ} in Dhao) in Bima
and \tsc{Havu-Dhao} are given in \trf{03-16}
and \trf{03-17}, along with reflexes from Sanggar
and Kambera to show that this change is exclusive to Bima and \tsc{Havu-Dhao}.
In addition to these examples,
*qatay ‘liver’ (\trf{03-11})
and *batu ‘stone’ (\trf{03-01} also share *t > \it{d} in Bima and \tsc{Havu-Dhao}.

\begin{table}[p]
	\centering\tcp{Bima and \tsc{Havu-Dhao} *t > **d /{\#\gap}}{03-16}
	\st{4pt}\begin{tabular}{ll@{}llllll}\lsptoprule
*gloss&afraid&rope&indigo&person&sugarcane&mat&eel\\
PMP&*(ma)-\tbr{t}akut&*\tbr{t}alih&*\tbr{t}aRum&*\tbr{t}au&*\tbr{t}əbuh&*\tbr{t}əpiR&*\tbr{t}una\\ \midrule
Bima&{\hr}\tbr{d}ahu&{\hr}\tbr{d}ari&{\hr}\tbr{d}au&{\hr}\tbr{d}ou&{\hr}\tbr{d}oɓu&{\hr}\tbr{d}ipi&{\hr}\tbr{d}una\\
Havu&{\hr}me\tbr{d}aʔu&{\hr}\tbr{d}ari&{\hr}\tbr{d}ao&{\hr}\tbr{d}əu&{\hr}\tbr{ɗ}əbu&{\hr}\tbr{d}əpi&{\hr}\tbr{d}əno\\
Dhao&{\hr}ma\textcolor{DarkRed}{\tbfDZ}aʔu&{\hr}\textcolor{DarkRed}{\tbfDZ}ari&&{\hr}\textcolor{DarkRed}{\tbfDZ}əu&{\hr}\textcolor{DarkRed}{\tbfDZ}əbu&{\hr}\textcolor{DarkRed}{\tbfDZ}əpi&\cg{(roma)}\\ \midrule
Sanggar&\cg{(‹podo›)}&&‹taogg›&\cg{(ata)}&{\hr}tewu&&\\
Kambera&\cg{(maŋɐɗat-u)}&\cg{(liku)}&\cg{(wora)}&{\hr}tau&{\hr}tiɓu&{\hr}topu&{\hr}tuna\\
		\lspbottomrule
	\end{tabular}
\end{table}

\begin{table}[p]\centering\tcp{Bima and \tsc{Havu-Dhao} *t > **d /V{\gap}V}{03-17}
\begin{adjustbox}{width=\textwidth}
	\begin{threeparttable}\st{2pt}
	\begin{tabular}{llllllllll}\lsptoprule
*gloss&blind&see&louse&raw&eye&die&seven&egg&cucumber\\
PMP&*bu\tbr{t}a&*ki\tbr{t}a&*ku\tbr{t}u&*ma-ha\tbr{t}aq&*ma\tbr{t}a&*ma\tbr{t}ay&*pi\tbr{t}u&*qa\tbr{t}əluR&*qa\tbr{t}imun\\ \midrule
Bima&{\hr}mbu\tbr{d}a&{\hr}e\tbr{d}a&{\hr}hu\tbr{d}u&{\hr}ma\tbr{d}a&{\hr}ma\tbr{d}a&{\hr}ma\tbr{d}e&{\hr}pi\tbr{d}u&{\hr}\tbr{d}olu&{\hr}\tbr{d}imu\\
Havu&{\hr}ɓə\tbr{d}u&{\hr}ŋə\tbr{d}i&{\hr}u\tbr{d}u&{\hr}ma\tbr{d}a&{\hr}ma\tbr{d}a&{\hr}ma\tbr{d}e&{\hr}pi\tbr{d}u&{\hr}\tbr{d}əlu&{\hr}vo\tbr{d}imu\su{†}\\
Dhao&{\hr}bə\textcolor{DarkRed}{\tbfDZ}u&{\hr}-ə\textcolor{DarkRed}{\tbfDZ}i&{\hr}u\textcolor{DarkRed}{\tbfDZ}u&{\hr}mama\textcolor{DarkRed}{\tbfDZ}a&{\hr}ma\textcolor{DarkRed}{\tbfDZ}a&{\hr}ma\textcolor{DarkRed}{\tbfDZ}e&{\hr}pi\textcolor{DarkRed}{\tbfDZ}u&\cg{(kanaɖ͡ʐu)}&\\ \midrule
Sanggar&\cg{(piso)}&&{\hr}kutu&&{\hr}mata&{\hr}mate&{\hr}pitu&{\hr}telu&\\
Kambera&\cg{(poki)}&{\hr}ita&{\hr}wutu&&{\hr}mata&{\hr}meti&{\hr}pihu&{\hr}tilu&{\hr}kaɗimbu\\
		\lspbottomrule
	\end{tabular}
		\begin{tablenotes}\tnsize
			\item[†]	{Havu has \it{vodimu poro} ‘cucumber’ and \it{vodimu dana} ‘watermelon’ \citep[18]{Wijngaarden1896}.}
		\end{tablenotes}
	\end{threeparttable}
\end{adjustbox}
\end{table}

\begin{table}[p]\centering\tcp{Bima *t = \it{t}}{03-18}
	\begin{threeparttable}\st{3.5pt}
	\begin{tabular}{llll|lll}\lsptoprule
*gloss&faeces&palm wine&we \cg{(incl.)}&low tide\su{‡}&three&knee\\
PMP&*\tbr{t}aqi&*\tbr{t}uak&*ki\tbr{t}a&*ma-qəti&*\tbr{t}əlu&*\tbr{t}uqud\\	\midrule
Bima&{\hr}\tbr{t}aʔi&{\hr}\tbr{t}ua&{\hr}ndai\tbr{t}a&{\hr}mo\tbr{t}i&{\hr}\tbr{t}olu&{\hr}ta\tbr{t}uʔu\\
Havu&‹\tbr{d}ei›&{\hr}\tbr{d}ue&{\hr}\tbr{d}i&{\hr}mə\tbr{t}i&{\hr}\tbr{t}əlu&{\hr}ru\tbr{t}uu\\
Dhao&{\hr}\textcolor{DarkRed}{\tbfDZ}əi&{\hr}nau-\textcolor{DarkRed}{\tbfDZ}ua\su{†}&{\hr}ə\textcolor{DarkRed}{\tbfDZ}i&{\hr}mə\tbr{t}i&{\hr}\tbr{t}əlu&{\hr}uru\tbr{t}uu\\	\midrule
Sanggar&\cg{(poʤa)}&&&{\hr}meti&{\hr}telu&{\hr}tuu\\
Kambera&{\hr}tai&&{\hr}ɲuta&{\hr}moti&{\hr}tilu, tailu&\cg{(kambɐku)}\\
		\lspbottomrule
	\end{tabular}
		\begin{tablenotes}\tnsize
			\item[†]	{Dhao \it{nau-ɖ͡ʐua} is ‘cluster of lontar fruit’.}
			\item[‡]	{Bima \it{moti} is ‘sea’.
								Havu and Dhao \it{məti} are ‘dry’.}
		\end{tablenotes}
	\end{threeparttable}
\end{table}

The six words which retain *t = \it{t} in Bima for which a cognate
has been identified in \tsc{Havu-Dhao} are given in \trf{03-18}.
Of these, three also retain *t = \it{t} in \tsc{Havu-Dhao},
and three have *t > \it{d} in \tsc{Havu-Dhao}.
The single word in which Bima has *t > \it{d}
and \tsc{Havu-Dhao} have *t = \it{t} is *təktək ‘tokay gecko’
> Bima \it{deke}, Havu \it{teke}, Dhao \it{təke}.
The onomatopoeic nature of this word may account for the irregularity.

Examples in which Bima has *p > \it{f}
and \tsc{Havu-Dhao} has *p > {\0} are given in \trf{03-19},
and examples of *p = \it{p} are given in \trf{03-20}.
\citet[93]{Blust2008} observes that *p = \it{p} is the usual
reflex after schwa or in historic nasal-plosive clusters.
Two additional examples which show *p = \it{p} are *təpiR ‘mat’
(\trf{03-16}) and *pitu ‘seven’ (\trf{03-17}).

\begin{table}\tcp{Bima and \tsc{Havu-Dhao} *p > **f}{03-19}
\begin{adjustbox}{width=\textwidth}
	\begin{threeparttable}\st{2pt}
	\begin{tabular}{llllllll}\lsptoprule
*gloss&rice plant&stingray&turtle&dream&fire&centipede&calcium\\
PMP&*\tbr{p}ajay&*\tbr{p}aRih&*\tbr{p}əñu&*h\<in\>i\tbr{p}i&*ha\tbr{p}uy&*qalu-hi\tbr{p}an&*qa\tbr{p}uR\\	\midrule
Bima&{\hr}\tbr{f}are&{\hr}\tbr{f}ai&{\hr}\tbr{f}onu&{\hr}ni\tbr{f}i&{\hr}a\tbr{f}i&{\hr}ri\tbr{f}a&{\hr}a\tbr{f}u\\
Havu&\hp{*\tbr{f}}are&&\hp{*\tbr{f}}əɲu&{\hr}nii&{\hr}ai&{\hr}terie rai&{\hr}ao\\
Dhao&\hp{*\tbr{f}}are&\hp{*\tbr{f}}ai&\hp{*\tbr{f}}əɲu&{\hr}nii&{\hr}ai&&{\hr}ao\\ \midrule
Sanggar&{\hr}pare&&&‹ngenipekh›&{\hr}api&&{\hr}apu\\
Kambera&{\hr}pari&{\hr}pai, pari&\cg{(tanoma)}\su{†}&\cg{(maŋgaɗipaŋ-u)}\su{‡}&{\hr}epi&{\hr}kalipa&{\hr}kapu\\
		\lspbottomrule
	\end{tabular}
		\begin{tablenotes}\tnsize
			%\item[†]	{Kambera \it{pari} is ‘rice’ \citep[404]{Onvlee1984}.
								%It has \it{uhu} for ‘rice plant’.}
			\item[†]	{Weyewa has \it{wonːu} and Kodi has \it{ɣanu} ‘turtle’ \citep[472]{Onvlee1984}.}
			\item[‡]	{Weyewa has \it{kanaŋganipi} ‘dream’ \citep[271]{Onvlee1984}
									which contains a reflex of *h\<in\>ipi.}
		\end{tablenotes}
	\end{threeparttable}
\end{adjustbox}
\end{table}

\begin{table}
	\centering\tcp{Bima and \tsc{Havu-Dhao} *p = \it{p}}{03-20}
	\begin{adjustbox}{width=\textwidth}
	\st{2pt}\begin{tabular}{lllllllll}\lsptoprule
*gloss&fart&how m.?&hot&fathom&four&sick&fold&PP/CC\\
PMP&**\tbr{p}əsu&*\tbr{p}ija&*(ma)-\tbr{p}anas&*\tbr{d}əpa&*ə\tbr{p}at&*ha\tbr{p}əjəs&*lə\tbr{p}ət&*ə\tbr{mp}u\\	\midrule
Bima&{\hr}\tbr{p}oʧu&{\hr}\tbr{p}ila&{\hr}\tbr{p}ana, \tbr{mb}ana&{\hr}ndu\tbr{p}a&{\hr}u\tbr{p}a&{\hr}\tbr{p}ili&{\hr}li\tbr{p}i&{\hr}o\tbr{mp}u\\
Havu&{\hr}\tbr{p}əhu&{\hr}\tbr{p}əri&{\hr}\tbr{p}ana&{\hr}rə\tbr{p}a&{\hr}ə\tbr{p}a&{\hr}\tbr{p}əɗa&{\hr}lə\tbr{p}a&{\hr}ə\tbr{p}u\\
Dhao&{\hr}\tbr{p}əʧu&{\hr}\tbr{p}əri&{\hr}\tbr{p}ana&{\hr}rə\tbr{p}a&{\hr}ə\tbr{p}a&{\hr}\tbr{p}əda&{\hr}lə\tbr{p}e&{\hr}ə\tbr{p}u\\	\midrule
Sanggar&{\hr}pesu&{\hr}pires&{\hr}pana&{\hr}ndupa&&&&\\
Kambera&{\hr}kapihu&{\hr}piraŋ-u&{\hr}mbana&{\hr}rɐpa, dɐpa&{\hr}patu&\cg{(hiɗu)}&{\hr}lɐpit&{\hr}umbuk\\
		\lspbottomrule
	\end{tabular}
	\end{adjustbox}
\end{table}

The two examples in which Bima has *p = \it{p}
and \tsc{Havu-Dhao} has *p > {\0} are *paniŋ ‘feed’ > Bima \it{pani},
Havu, Dhao \it{ani}, and *ma-nipis ‘thin’ > Bima \it{nipi},
Havu \it{menii}, Dhao \it{manii}.
There is also one example which shows a split in Havu: *qapəju
‘gall’ > Bima \it{folu},
Havu \it{əɗu} ‘gall’ alongside *qapəju > Havu \it{pəɗu}, Dhao \it{papədu} ‘bitter’.
The \tsc{Havu-Dhao} words meaning ‘bitter’ are probably from
intermediate *ma-pəju > **mpəju,
thus attesting a historic nasal-plosive cluster.
There is only a single example in which Bima has *p > \it{f}
and \tsc{Havu-Dhao} has *p = \it{p}; *pahuq ‘mango’ > Bima \it{foʔo},
Havu, Dhao \it{pau}.\footnote{
		One word with PMP *p which has
		a cognate in both Bima
		and Havu-Dhao has an irregular reflex:
		*pənuq ‘full’ > Bima \it{ɓini}, Dhao \it{pənu}.
		Havu has non-cognate \it{tobo} ‘full’.}

\clearpage
Examples of shared *l > \it{r} in Bima
and \tsc{Havu-Dhao} are given in \trf{03-21}.
In addition, *balabaw ‘mouse, rat’ (\trf{03-12}), *laRiw ‘run (\trf{03-09}),
*talih ‘rope’ (\trf{03-16}),
and *qalu-hipan ‘centipede’ (\trf{03-19}) also share *l > \it{r}.

\begin{table}\tcp{Bima and \tsc{Havu-Dhao} *l > \it{r}}{03-21}
\begin{adjustbox}{width=\textwidth}
	\begin{threeparttable}\st{2pt}
	\begin{tabular}{llllllll}\lsptoprule
*gloss&moon&grain head&body hair&skin&head\su{‡}&ten&eight\\
PMP&*bu\tbr{l}an&*bu\tbr{l}iR&*bu\tbr{l}u&*ku\tbr{l}it&*qu\tbr{l}u&*sa-ŋa-pu\tbr{l}uq&*wa\tbr{l}u\\	\midrule
Bima&{\hr}wu\tbr{r}a&{\hr}wu\tbr{r}i&{\hr}wu\tbr{r}u&{\hr}hu\tbr{r}i&{\hr}u\tbr{r}u&{\hr}sampu\tbr{r}u&{\hr}wa\tbr{r}u\\
Havu&{\hr}və\tbr{r}u&{\hr}vu\tbr{r}i&{\hr}vu\tbr{r}u\su{†}&{\hr}ku\tbr{r}i&{\hr}u\tbr{r}u&{\hr}heŋu\tbr{r}u&{\hr}a\tbr{r}u\\
Dhao&{\hr}hə\tbr{r}u&{\hr}hu\tbr{r}i&&{\hr}kaʔu\tbr{r}i&{\hr}u\tbr{r}u&{\hr}ʧaŋu\tbr{r}u&{\hr}a\tbr{r}u\\	\midrule
Sanggar&‹woelokh›&&&{\hr}loki&&&{\hr}walu\\
Kambera&{\hr}wulaŋ-u&{\hr}wili&{\hr}wulu&{\hr}kalit-u&&{\hr}hakambulu&{\hr}walu\\
		\lspbottomrule
	\end{tabular}
		\begin{tablenotes}\tnsize
			\item[†]	{Havu \it{vuru} is ‘fibres, root fibres’ \citep[119]{Wijngaarden1896}.}
			\item[‡]	{Bima \it{uru} is ‘headwaters’,
								Havu and Dhao \it{uru} is ‘formerly’.}
		\end{tablenotes}
	\end{threeparttable}
\end{adjustbox}
\end{table}

Examples in which Bima
and Havu-Dhao retain *l = \it{l} are given in \trf{03-22}.
In addition to these,
*l = \it{l} is also found in *ləpət ‘fold’ (\trf{03-20}),
*qaləjaw ‘sun, day’ (\trf{03-12}), and *qatəluR ‘egg’ (\trf{03-17}).

\begin{table}
	\centering\tcp{Bima and \tsc{Havu-Dhao} *l = \it{l}}{03-22}
\begin{adjustbox}{width=\textwidth}
	\st{2pt}\begin{tabular}{lllllllll}\lsptoprule
*gloss&go&sail&wax, candle&roll up&five&price&three&wide\\
PMP&*\tbr{l}akaw&*\tbr{l}ayaR&*\tbr{l}i\tbr{l}in&*\tbr{l}ulun&*\tbr{l}ima&*bə\tbr{l}i&*tə\tbr{l}u&*bəkə\tbr{l}aj\\	\midrule
Bima&{\hr}\tbr{l}ao&{\hr}\tbr{l}oʤa&{\hr}\tbr{l}i\tbr{l}i&{\hr}\tbr{l}uru&{\hr}\tbr{l}ima&{\hr}we\tbr{l}i&{\hr}to\tbr{l}u&{\hr}we\tbr{l}a\\
Havu&{\hr}\tbr{l}a&{\hr}\tbr{l}ai&{\hr}\tbr{l}i\tbr{l}i&{\hr}\tbr{l}ule&{\hr}\tbr{l}əmi&{\hr}və\tbr{l}i&{\hr}tə\tbr{l}u&{\hr}bə\tbr{l}a, keɓə\tbr{l}a\\
Dhao&{\hr}\tbr{l}a-&{\hr}\tbr{l}ai&{\hr}\tbr{l}i\tbr{l}i&{\hr}\tbr{l}ulu&{\hr}\tbr{l}əmi&{\hr}hə\tbr{l}i&{\hr}tə\tbr{l}u&{\hr}ɓə\tbr{l}a\\	\midrule
Sanggar&{\hr}lako&{\hr}loʤa&‹lilikh›&&{\hr}lima&{\hr}weli&{\hr}telu&\\
Kambera&{\hr}laku&{\hr}liru&{\hr}liliŋ-u&\cg{(lukuru)}&{\hr}lima&{\hr}wili&{\hr}tilu, tailu&{\hr}mbelar-u\\
		\lspbottomrule
	\end{tabular}
\end{adjustbox}
\end{table}

Examples in which Bima has *l > \it{r} but Havu-Dhao retain *l
= \it{l} are given in \trf{03-23}.
Two cognate sets show a split of *l > \it{r,
l} in individual languages.
Firstly, *qahəlu ‘pestle’ > Bima \it{aru},
Dhao \it{aru} but Havu \it{aru {\tl} alu}.
Secondly, PMP *balik ‘reverse, turn around’ and/or *baliw ‘return’
> Bima \it{mbali} ‘turn’,
\it{wali} ‘again’, \citep[161]{Ismail1985} \it{wari} ‘turn’ \citep[181]{Zollinger1850},
Havu \it{ɓari} ‘turn’, \it{keɓəle/keɓəli} ‘turned’, Dhao \it{bari/kabəli} ‘turn’.

\begin{table}\tcp{Bima *l > \it{r}}{03-23}
	\begin{threeparttable}\st{6pt}
	\begin{tabular}{llllll}\lsptoprule
*gloss&widow, widower&male\su{†}&ginger&roll up\\
PMP&*ba\tbr{l}u&*\tbr{l}aki&*\tbr{l}aqia&*lu\tbr{l}un\\ \midrule
Bima&{\hr}mba\tbr{r}u \it{‘bachelor’}&{\hr}\tbr{r}ahi&{\hr}\tbr{r}ea&{\hr}lu\tbr{r}u\\
Havu&{\hr}ɓa\tbr{l}u&{\hr}\tbr{l}aʔi&&{\hr}lu\tbr{l}e\\
Dhao&{\hr}ba\tbr{l}u&{\hr}\tbr{l}aʔi&{\hr}\tbr{l}ia pana&{\hr}lu\tbr{l}u\\ \midrule
Sanggar&{\hr}walu mone&&&\\
Kambera&\cg{(londaŋu)}&{\hr}lei&{\hr}layia&\cg{(lukuru)}\\
		\lspbottomrule
	\end{tabular}
		\begin{tablenotes}\tnsize
			%\item[†]	{Bima \it{mbaru} is ‘bachelor’.}
			\item[†]	{Bima \it{rahi} is ‘husband’,
								Havu \it{laʔi} is ‘husband,
								wife’, and Dhao \it{laʔi} is ‘male’.}
		\end{tablenotes}
	\end{threeparttable}
\end{table}

Among the splits potentially shared between Bima
and Havu-Dhao, the split of *s > \it{ʧ,
s} provides the weakest evidence of being shared,
as this change is only unambiguously attested in Bima
and Dhao, as Havu has *s > \it{h} in all cases.
Examples of *s = \it{s}
and *s > \it{ʧ} in Bima are given in \trf{03-24}
and \trf{03-25} respectively.
Manggarai reflexes are also given for comparison as it has regular *s > \it{ʧ}.
Note, however, that other languages of Flores,
with which Manggarai is more closely related,
do not have *s > \it{ʧ}.
Thus, whatever the status of shared *s > \it{ʧ} in Bima
and Dhao, the change in Manggarai may be due to parallel development.


\begin{table}\tcp{Bima and Dhao *s = \it{s}}{03-24}
	\begin{threeparttable}\st{6pt}
	\begin{tabular}{lllllll}\lsptoprule
*gloss&oar&contents&alone&breast&salt&island\\
PMP&*bəR\tbr{s}ay&*i\tbr{s}i&*ma-ə\tbr{s}a&*\tbr{s}u\tbr{s}u&*qa\tbr{s}iRa&*nu\tbr{s}a\\	\midrule
Bima&{\hr}we\tbr{s}e&{\hr}i\tbr{s}i&{\hr}ke\tbr{s}e&{\hr}\tbr{s}u\tbr{s}u&{\hr}\tbr{s}ia&{\hr}ni\tbr{s}a\\
Dhao&{\hr}\tbr{s}ehe\su{†}&{\hr}i\tbr{s}i&{\hr}me\tbr{s}a&{\hr}\tbr{s}u\tbr{s}u&\cg{(həro)}&\cg{(kabarai)}\\
Havu&\cg{(tuku)}&{\hr}i\tbr{h}i&{\hr}mi\tbr{h}a&{\hr}\tbr{h}u\tbr{h}u&\cg{(masi)}&\\	\midrule
Sanggar&\cg{(krawe)}&{\hr}panisi&{\hr}mesa&{\hr}susu&{\hr}sia&\\
Kambera&{\hr}buhi&{\hr}ihi&{\hr}meha&{\hr}huhu&\cg{(mehi)}&{\hr}nuha\\
Manggarai&\cg{(bise)}\su{†}&{\hr}iʧi&\cg{(hanaŋ)}&{\hr}ʧuʧu&{\hr}ʧiʔe&{\hr}nuʧa\\
		\lspbottomrule
	\end{tabular}
		\begin{tablenotes}\tnsize
			\item[†]	{Dhao \it{sehe} has metathesis of the initial and medial consonants.
								It is probably a borrowing from a Rote language,
								all of which have \it{sefe} ‘oar’ with the same metathesis.
								Manggarai \it{bise} ‘paddle’ is given as a borrowing
								from Makasar by \citet[vol. 2; 527]{Verheijen1995}.}
		\end{tablenotes}
	\end{threeparttable}
\end{table}

\begin{table}\tcp{Bima *s > \it{ʧ}}{03-25}
\begin{adjustbox}{width=\textwidth}
	\begin{threeparttable}\st{2pt}
	\begin{tabular}{llllllllll}\lsptoprule
*gloss&who?&nine&one&mortar&fart&branch&wet&elbow&salty\\
PMP&*\tbr{s}ai&*siwa&*isa&*ləsuŋ&**pəsu&*saŋa&*basəq&*siku&*ma-qasin\\
Bima&{\hr}\tbr{ʧ}ou{\cgr}&{\cgr}{\hr}\tbr{ʧ}iwi{\cgr}&{\hr}i\tbr{ʧ}a{\cgr}&{\hr}no\tbr{ʧ}u{\cgr}&{\hr\hr}po\tbr{ʧ}u{\cgr}&{\hr}\tbr{s}aŋa,{\cgr}&{\hr}mbe\tbr{ʧ}a&{\hr}\tbr{ʧ}ihu&{\hr}ma\tbr{ʧ}i\su{†}\\
&{\cgr}&{\cgr}&{\cgr}&{\cgr}&{\cgr}&{\hr}\tbr{ʧ}aŋa\su{‡}{\cgr}&&& \\
Dhao&{\hr}\tbr{ʧ}ee{\cgr}&{\hr}\tbr{ʧ}eo{\cgr}&{\hr}ə\tbr{ʧ}i{\cgr}&{\hr}ŋə\tbr{ʧ}u{\cgr}&{\hr\hr}pə\tbr{ʧ}u{\cgr}&{\hr}ka\tbr{ʧ}aŋa{\cgr}&{\hr}ba\tbr{s}a&{\hr}hui \tbr{s}iʔu&{\hr}ma\tbr{s}i\\
Havu&\cg{(naduu)}&{\hr}\tbr{h}eo&{\hr}ə\tbr{h}i&{\hr}ŋə\tbr{h}u&{\hr\hr}pə\tbr{h}u&{\hr}\tbr{ʤ}aŋa&{\hr}ɓa\tbr{h}a&{\hr}vo\tbr{h}iu&\cg{(həro)}\\ \midrule
Sanggar&{\hr}sei&{\hr}siwa&&&{\hr\hr}pesu&{\hr}ʧaŋat&\cg{(ŋgesak)}&{\hr}siku&\\
Kambera&{\hr}ŋgaa&{\hr}hiwa&{\hr}ɗiha&{\hr}ŋohu&{\hr\hr}kapihu&{\hr}ɲʤaŋa&{\hr}mbaha&{\hr}hiu&{\hr}mehi\\
Manggarai&{\hr}ʧei&{\hr}ʧiok&{\hr}ʧaa&{\hr}ŋəɲʧuŋ&{\hr\hr}pəʧu&\cg{(daŋka)}&{\hr}baʧa&{\hr}ʧiku&\cg{(ʧərak)}\\
		\lspbottomrule
	\end{tabular}
		\begin{tablenotes}\tnsize
			\item[†]	{Bima \it{maʧi} is ‘sweet’.}
			\item[‡]	{\citet[139]{Ismail1985} \it{saŋa} give ‘forked branch,
								branch’ \citet[206]{Zollinger1850} gives ‹tjanga› ‘forked branch’.}
		\end{tablenotes}
	\end{threeparttable}
\end{adjustbox}
\end{table}
